%%%%%%%%%%%%%%%%%%%%%%%%%%%%%%%%%%%%%%%%%%%%%%%%%%%%%%%%%%%
%%%
%%%            日本伝熱学会主催「日本伝熱シンポジウム」
%%%  【非公式】講演アブストラクト原稿（A4 用紙 2--6 枚，日本語）
%%%
%%%     <https://github.com/Yuki-MATSUKAWA/NHTS_template>
%%%
%%%                      v1.0.0 Yuki MATSUKAWA xx Apr. 2024
%%%
%%%%%%%%%%%%%%%%%%%%%%%%%%%%%%%%%%%%%%%%%%%%%%%%%%%%%%%%%%%

\documentclass[
    paper=a4paper,      % A4 用紙サイズ
    article,            % article 相当の文書クラス
    fleqn,              % 数式を左寄せ
    fontsize=10pt,      % 欧文サイズ 10 pt
    jafontsize=10pt,    % 和文サイズ 10 pt
    head_space=20mm,    % 天の余白
    foot_space=25mm,    % 地の余白
    gutter=20mm,        % のどの余白
    fore-edge=20mm      % 小口の余白
    ]{jlreq}            % jlreq クラスを使用

%%% 使用するスタイルファイル %%%
\usepackage{settings}

%%% フッターの情報 %%%
\newcommand{\foot}{
  %%% 開催する回に応じて内容を忘れずに変更 %%%
  第xx回日本伝熱シンポジウム講演論文集 (20xx-5)
  %%%%%%%%%%%%%%%%%%%%%%%%%%%%%%%%%%%%%%%%%
  }

\begin{document}

\thispagestyle{nhts}

%%% 著者情報 %%%
\begin{center}
    \fontsize{14pt}{18pt}\selectfont
    {\gtfamily
    %%%%%%%%%%%%%%%%%%%%%%%%%%%%%%%
    ここには和文表題を書いてください．
    %%%%%%%%%%%%%%%%%%%%%%%%%%%%%%%
    }\\\fontsize{12pt}{14pt}\selectfont
    {\gtfamily
    %%%%%%%%%%%%%%%%%%%%%%%%%%%%%%%
    ここには和文副題を書いてください．
    %%%%%%%%%%%%%%%%%%%%%%%%%%%%%%%
    }\\
    %%%%%%%%%%%%%%%%%%%%%%%%%%%%%%%%%%%%
    Please write the English title here.
    %%%%%%%%%%%%%%%%%%%%%%%%%%%%%%%%%%%%
    \\
    %%%%%%%%%%%%%%%%%%%%%%%%%%%%%%%%%%%%%%%
    Please write the English subtitle here.
    %%%%%%%%%%%%%%%%%%%%%%%%%%%%%%%%%%%%%%%
    \vskip12pt\fontsize{10pt}{12pt}\selectfont
    \begin{tabular}{lrll}
      伝学  &$^\ast$\hspace{-3mm} &伝熱 太郎  &（伝熱大学）\\
        & &伝熱 花子  &（伝熱大学）\\
      伝正  & &抜山 四郎  &（東京伝熱大学）
    \end{tabular}
    \vskip12pt
    Taro DENNETSU$^1$, Hanako DENNETSU$^1$ and Shiro NUKIYAMA$^2$
    \\
    $^1$Dept. of Mech. Eng., Dennetsu Univ., 5-1-5, Kashiwanoha, Kashiwa, 277-8563 \\
    $^2$Dept. of Mech. Eng., Tokyo Dennetsu Univ., 1-18-11, Uchikanda, Chiyoda, 101-0047
    \vskip12pt
\end{center}
{\setlength{\leftskip}{\leftskip+15truemm}\setlength{\rightskip}{\rightskip+15truemm}%
%%% ここにはアブストラクトを書いてください %%%
\lipsum[1-2]
%%%%%%%%%%%%%%%%%%%%%%%%%%%%%%%%%%%%%%%%%%%
\par}%
\begin{center}
    \vskip12pt
    \textit{Key Words}: 
    %%% ここにはこの研究のキーワードを書いてください %%%
    Heat Transfer, Boiling, Two-phase flow, Forced convection, Natural convection
    %%%%%%%%%%%%%%%%%%%%%%%%%%%%%%%%%%%%%%%%%%%%%%%%
    \vskip12pt
\end{center}

\section{序論}

\jalipsum[1-4]{wagahai}

\lipsum[1-4]

\section{解析対象・計算条件}

\jalipsum[1-4]{wagahai}

\subsection{支配方程式}

\subsection{計算条件}

\begin{table}[tp]
  \centering
  \caption{Sample table.}
  \label{tab:table}
  \begin{tabularx}{0.5\textwidth}{C|C|C} \hline\hline
    A  &a  &1 \\ \hline
    B  &b  &2 \\ \hline
    C  &c  &3 \\ \hline
    D  &d  &4 \\ \hline
  \end{tabularx}
\end{table}

\subsubsection{計算コードの妥当性検証}

\lipsum[1-4]

\jalipsum[1-4]{wagahai}

\section{結果・考察}

\begin{figure}[tp]
  \centering
  \includegraphics[width=0.5\columnwidth]{figure/tiger.pdf}
  \caption{A sample figure.}
  \label{fig:one_figure}
\end{figure}


\begin{figure}[tp]
  \centering
  \begin{subfigure}{0.45\columnwidth}
      \centering
      \includegraphics[width=\columnwidth]{example-image-a}
      \caption{Left figure.}
      \label{fig:two_figures_a}
  \end{subfigure}
  \hfill % ここで空白を入れると図が適切に配置される．\hspace{} なども可．
  \begin{subfigure}{0.45\columnwidth}
      \centering
      \includegraphics[width=\columnwidth]{example-image-b}
      \caption{Right figure.}
      \label{fig:two_figures_b}
  \end{subfigure}
  \caption{Two sample figures.}
  \label{fig:two_figures}
\end{figure}

\lipsum[1-4]

\section{結論}

\jalipsum[1-4]{wagahai}

\section*{謝辞}

\lipsum[1-4]

\end{document}
